\section{The last partially filled layer and the nearest zero}
Here \(k=k(n)\to\infty\), \(k^2 8^k=o(n)\), \(N=2^n\), and \(C_k\) denotes the minimum gate count of a uniformly chosen truth table of weight \(k\). Put
\[
 \mu=n2^{1-k},\qquad \sigma^2=\mu(1-2^{1-k}).
\]
For these assertions it suffices to have the sparse-support bounds \(2n-Z-O(k+\log k)\le C\le2n-Z+O(k2^k)\), where \(Z\) counts constant raw columns of the \(k\) support addresses. The full binary basis and counted De Morgan basis obey these inequalities by the sparse-support lower bound and repeated-column construction used above. Red'kin's \href{https://www.mathnet.ru/eng/dm172}{2004 theorem} analyzes the Finikov synthesis method over the binary basis with XOR and XNOR excluded.

For independent random rows, \(Z\) is binomial with mean \(\mu\). Conditioning them to be distinct has total-variation cost at most \(\binom{k}{2}/N\). The binomial Berry--Esseen inequality, applied to \(-Z+\mu\), and the two deterministic shifts of size \(O(k2^k)\) therefore give
\[
 \sup_x\left|\Pr\!\left[\frac{C_k-(2n-\mu)}{\sigma}\le x\right]
        -\Phi_{\rm G}(x)\right|
 \ \le O\!\left(\frac{k2^k+1}{\sigma}+\frac{k^2}{N}\right)=o(1),
\]
where \(\Phi_{\rm G}\) is the standard normal distribution function. Translating the Gaussian argument costs at most the shift times its bounded density.

Take \(s=\lfloor2n-\mu+x_n\sigma\rfloor\), with \(|x_n|=o(\sqrt\mu)\), and let \(\lambda_k=\Pr[C_k\le s]\) be the \emph{actual} accepted fraction on layer \(k\). Write \(r=W(1/e)\), \(b=r/(1+r)\), \(\lambda_*=1/(1+r)\) and \(a=r/(1+r)^2\). Let \(R_N\) be the least modulus of a zero of \(P_0\) at this budget.
\begin{theorem}
If \(\sqrt{k}\,|\lambda_k-\lambda_*|\to\infty\), then
\[
 NR_N=rk+b\left[\tfrac12\log(2\pi k)-\log|\lambda_k-\lambda_*|\right]+o(1).
\]
Every nearest zero \(\zeta\) satisfies \(\sqrt{k}\,|\arg(-\zeta)|\to0\) in this case; \(\arg\) is the principal argument. If \(\sqrt{k}(\lambda_k-\lambda_*)\to h\in\mathbb R\), put
\[
 B_h(t)=\frac{r t^2}{2}-\frac12\log(h^2+a^2t^2),\qquad
 t_h=\sqrt{\max(1/r-h^2/a^2,0)} ,
\]
interpreting a zero logarithm argument as \(B_h(t)=+\infty\). Then
\[
 NR_N=rk+b\left[\log k+\tfrac12\log(2\pi)
                       +B_h(t_h)\right]+o(1),\qquad
       \sqrt{k}\,|\arg(-\zeta)|\longrightarrow t_h .
\]
For instance the Gaussian approximation can be used in place of \(\lambda_k\) at the critical precision if \(k^3 8^k=o(n)\).
\end{theorem}

\begin{proof}
Let \(E_j(u)=\sum_{i=0}^j u^i/i!\) and \(G_k(w)=P_0(kw/N)\). On each fixed disk \(|w|\le R\), for every fixed \(A>0\),
\[
  G_k(w)=E_{k-1}(kw)+\lambda_k(kw)^k/k!+O_{A,R}(e^{-Ak}).
\]
Indeed the constant-column mean is at least \(2\mu\) below layer \(k\) and at most \(\mu/2\) above it. In any range \(1\le j\le Dk\), for fixed \(D\), the lower-layer upper synthesis error is \(O(k2^k)=o(\mu)\) and the upper-layer gate lower-bound error is \(O_D(k+\log k)=o(\mu)\). Chernoff bounds and distinctness conditioning thus make the incorrect fraction \(O(e^{-\Omega(\mu)}+(Dk)^2/N)\) outside layer \(k\). Layer zero is complete. The binomial-to-Poisson coefficient error through \(Dk\) is \(O_D(k^2/N)\) relative to the corresponding exponential-series coefficient before thinning. These errors are multiplied by at most \(e^{Rk}\). For the farther terms use the full exponential-series majorant and take \(D\) large enough that \(D(1+\log(R/D))<-A-1\).

On closed annuli \(0<\eta\le|w|\le R<1\) write exactly
\[
 G_k(w)=e^{kw}-\frac{(kw)^k}{k!}A_k(w),\qquad
 A_k(w)=\frac1{1-w}-\lambda_k+O_{\eta,R}(1/k).
\]
The same error estimate holds after one or two differentiations on smaller annuli. Indeed the normalized exponential tail is
\(\sum_{j\ge0}w^j\prod_{i=1}^j(1+i/k)^{-1}\), whose difference from \((1-w)^{-1}\) is bounded by \(\sum_{j\ge0}R^jj(j+1)/(2k)\). The preceding arbitrarily small exponential error remains negligible on the annulus after normalization. Cauchy's inequality gives the derivative bounds. The nonvanishing of \(e^{kw}\) and the same tail estimate exclude every disk \(|w|\le r-\eta\) from the zeros for large \(k\): there \(|w|e^{1+|w|}<1\) uniformly.

Write \(w=-\rho e^{i\theta}\) near \(-r\), \(d=k(\rho-r)\), and \(t=\sqrt k\,|\theta|\). For a zero, Stirling's formula gives the exact modulus equation up to a uniform \(O(k^{-1})\),
\[
 k(\Re w-1-\log\rho)=\log|A_k(w)|-\tfrac12\log(2\pi k)+O(k^{-1}).
\]
For \(|d|=O(\log k)\) and \(t=O(\sqrt{\log k})\) it becomes
\[
 b^{-1}d=\tfrac12\log(2\pi k)-\log|A_k(w)|+\tfrac12 rt^2+o(1).
 \tag{*}
\]
The error is uniform because \(k\theta^4+ d^2/k+|d|\theta^2=o(1)\) and \(-r-1=\log r\).

We first verify that the modulus equation really locates zeros to \(o(1)\) in \(d\). In the noncritical regime take \(\theta=0\); in the critical regime take \(\theta=t_h/\sqrt k\). In the latter case \(h^2+a^2t_h^2>0\). Uniformly through \(|d|\le C\log k\), the corresponding tail magnitude is asymptotic respectively to \(|\lambda_*-\lambda_k|\) or to \(k^{-1/2}\sqrt{h^2+a^2t_h^2}\): Taylor expansion of \((1-w)^{-1}\) gives the imaginary linear term \(-ia\theta\), with error \(O(\theta^2+|d|/k+1/k)\). Consequently the exact modulus equation has a solution \(w_0\) at the asserted value of \(d\) up to \(o(1)\), by the intermediate value theorem. More explicitly, insert into (*) the asserted expression for \(d\) plus and minus any fixed positive tolerance; the two residuals have opposite signs.

Choose analytic logarithms near \(w_0\) and write
\[
 T_k(w)=kw-k\log w-\log A_k(w)+\log k!-k\log k.
\]
The exact modulus equation is \(\Re T_k(w_0)=0\). A zero results whenever \(T_k(w)\in2\pi i\mathbb Z\). At \(w_0\), \(T'_k(w_0)=k(1+1/r)+O(\sqrt k+\log k)\), since \(|A_k(w_0)|\gtrsim k^{-1/2}\). Throughout a disk of radius \(C'/k\) for any fixed \(C'\) we have \(T''_k=O(k)\). Take \(2\pi i m\) nearest to \(T_k(w_0)\). The linearized root differs from \(w_0\) by \((2\pi i m-T_k(w_0))/T'_k(w_0)\), of magnitude \(O(1/k)\) and real part \(O(k^{-3/2})\). Taylor's theorem and Rouch\'e's theorem on a circle of radius \(C''/k^2\) about this linearized root, with sufficiently large \(C''\), give an exact root there. Its radial displacement from \(w_0\) is \(o(1/k)\), since \(\theta=O(k^{-1/2})\). This proves the stated expressions as upper bounds for \(NR_N\).

For a nearest zero these upper bounds give \(\rho\le r+O(\log k/k)\). The small-disk exclusion and the modulus equation, using \(\Re w\ge-\rho\) and boundedness of \(A_k\), now give \(d\ge-O(1)\) and \(k(1-\cos\theta)=O(\log k)\). Thus (*) applies to every nearest zero. Put \(H_k=\sqrt k\,(\lambda_k-\lambda_*)\). In this range
\[
 \sqrt k\,|A_k(w)|\le |H_k|+O(1+t).
\]
If \(|H_k|\to\infty\), comparison with the obtained upper bound shows \(rt^2/2\le\log(1+O(1+t)/|H_k|)+o(1)\). This first bounds \(t\) and then forces \(t\to0\), proving the matching radius by (*). For bounded \(H_k\), the same comparison yields \(rt^2/2\le\log(1+t)+O(1)\), so \(t\) is bounded again. A subsequential limit \(t'\) cannot satisfy \(h=t'=0\), since \(\sqrt k|A_k|\) would tend to zero in (*) in contradiction with the upper bound. Otherwise \(\sqrt k|A_k|\to\sqrt{h^2+a^2t'^2}\). Equation (*) now gives the correction \(B_h(t')\). Differentiation shows its unique minimum on \(t\ge0\) occurs at \(t_h\), also when \(h^2=a^2/r\). The upper bound forces this minimum for every subsequential limit, as claimed.
\end{proof}

For ordinary exponential sections, logarithmic corrections and the locations of individual zeros relative to the Szeg\H{o} curve have been studied by \href{https://doi.org/10.1216/rmjm/1181072998}{Carpenter, Varga and Waldvogel} and \href{https://www.math.kent.edu/~varga/pub/paper_178.pdf}{Varga and Carpenter}. Here the last coefficient can cancel the leading tail at the negative extremum, producing the critical scaling in the theorem.
